%!TEX root = ./main.tex

\section[Security]{The Arms Race}

% SECTION TIMING: 30:00--45:00 (7 frames). This section pays off the IMSI-catcher
% thread planted on Monday and hands directly to the O-RAN opener, which asks
% whether base-station security could be an app. Do not reorder those two.
%
% FRAMING BOUNDARY, enforced throughout: everything here describes capabilities and
% published research. Nothing here claims that anyone in the room is being intercepted.
% If a student asks "is this happening to me right now", the honest answer is that the
% equipment class is real and documented, and that we are not in a position to say more.

\begin{frame}{Monday's loose end}
  \vspace{0.3em}
  \centering
  \begin{tikzpicture}[scale=0.9,transform shape]
    \node[netnode,minimum width=1.8cm] (ph) at (0,0) {Phone};
    \node[corenode,minimum width=1.8cm] (nw) at (5.2,0) {Network};
    \draw[flow,draw=softgreen,very thick] (0.95,0.42) -- (4.25,0.42);
    \node[font=\small,anchor=south] at (2.6,0.5) {prove who you are};
    \draw[draw=coreorange,very thick,dashed] (4.25,-0.42) -- (0.95,-0.42);
    \node[font=\small,anchor=north] at (2.6,-0.5) {\textcolor{coreorange}{nothing comes back}};
  \end{tikzpicture}

  \vspace{0.9em}
  \small
  In 2G the phone proved itself to the network. The network proved nothing back,
  and the phone had already picked the loudest tower before anyone asked.

  \vspace{0.6em}
  \ask{Every generation since has been trying to close that gap. Did they?}
\end{frame}

\begin{frame}{The same problem, four times}
  \vspace{0.4em}
  {\small
  \begin{tabular}{@{}p{0.09\textwidth}p{0.40\textwidth}p{0.42\textwidth}@{}}
    \toprule
    & \textbf{\textcolor{coreorange}{The weakness}} & \textbf{\textcolor{softgreen}{The answer}}\\
    \midrule
    \textbf{2G} & Network never authenticates itself & ---\\[3pt]
    \textbf{3G} & Identity still sent in the clear at first contact & Mutual authentication added\\[3pt]
    \textbf{4G} & Same first-contact exposure & Stronger keys, same opening move\\[3pt]
    \textbf{5G} & Phone must still pick a tower first & Identity encrypted before it is sent\\
    \bottomrule
  \end{tabular}\par}

  \vspace{0.8em}
  \takeaway{Each generation closed a door. None of them closed the doorway.}

  \vspace{0.3em}
  {\scriptsize\textcolor{softgray}{Simplified for teaching. Each row is a family of
   changes across many releases, not a single switch that was flipped.}\par}
\end{frame}

\begin{frame}{Why the first message is the hard one}
  \begin{columns}[T,onlytextwidth]
    \column{0.52\textwidth}
      \centering
      \begin{tikzpicture}[scale=0.8,transform shape,
          step/.style={font=\scriptsize,align=left,text width=3.6cm}]
        \node[netnode,minimum width=1.5cm] (ph) at (0,0) {Phone};
        \node[oranode,minimum width=1.8cm] (bs) at (3.8,0) {Tower};
        \draw[flow,draw=ranpurple] (0.8,0.75) -- (2.9,0.75);
        \node[font=\scriptsize,anchor=south] at (1.85,0.8) {1. I want to attach};
        \draw[flow,draw=softgray] (2.9,0.05) -- (0.8,0.05);
        \node[font=\scriptsize,anchor=south] at (1.85,0.1) {2. Who are you?};
        \draw[flow,draw=coreorange,very thick] (0.8,-0.75) -- (2.9,-0.75);
        \node[font=\scriptsize,anchor=south] at (1.85,-0.7) {3. Here is my identity};
      \end{tikzpicture}

      \vspace{0.6em}
      {\scriptsize\textcolor{softgray}{Step 3 happens before the phone
       has verified anything about the tower.}\par}
    \column{0.44\textwidth}
      \large\textbf{The bootstrap problem.}

      \vspace{0.5em}
      \small
      To check a tower's credentials, the phone needs a channel.

      \vspace{0.5em}
      To get a channel, it has to talk to the tower.

      \vspace{0.5em}
      Something has to go first.
  \end{columns}

  \vspace{0.3em}
  \glossline{IMSI = International Mobile Subscriber Identity \quad
    SUPI = Subscription Permanent Identifier}
\end{frame}

\begin{frame}{5G stopped sending the identity in the clear}
  \begin{columns}[T,onlytextwidth]
    \column{0.50\textwidth}
      \centering
      \begin{tikzpicture}[scale=0.86,transform shape]
        \node[netnode,minimum width=1.9cm] (supi) at (0,0.9) {SUPI};
        \node[corenode,minimum width=2.4cm] (key) at (0,-0.15) {home network\\public key};
        \node[oranode,minimum width=1.9cm] (suci) at (0,-1.3) {SUCI};
        \draw[flow] (supi) -- (key);
        \draw[flow] (key) -- (suci);
        \node[font=\scriptsize,anchor=west,text width=2.9cm,align=left] at (1.5,-1.3)
          {this is what goes over the air};
      \end{tikzpicture}
    \column{0.46\textwidth}
      \large\textbf{Encrypt it first.}

      \vspace{0.5em}
      \small
      The permanent identity never leaves the phone unprotected.

      \vspace{0.5em}
      What the tower receives is a concealed version that only the home
      network can open.
  \end{columns}

  \glossline{SUCI = Subscription Concealed Identifier}

  \qa{So is the problem solved?}
     {Partly. It closes the easy collection of permanent identities --- but only
      where the operator has actually turned it on, and only if the phone stays
      on 5G.}
\end{frame}

\begin{frame}{Two ways it still leaks}
  \vspace{0.4em}
  \begin{columns}[T,onlytextwidth]
    \column{0.47\textwidth}
      \large\textbf{\textcolor{coreorange}{Configuration}}

      \vspace{0.5em}
      \small
      Concealment is a feature operators enable.

      \vspace{0.5em}
      A network that has not turned it on gets none of the benefit, and
      the phone cannot tell from the outside.
    \column{0.47\textwidth}
      \large\textbf{\textcolor{coreorange}{Downgrade}}

      \vspace{0.5em}
      \small
      Phones still support older generations so they keep working
      where 5G does not reach.

      \vspace{0.5em}
      Push a phone down to a weaker generation and you are back to that
      generation's rules.
  \end{columns}

  \vspace{0.8em}
  \takeaway{Backward compatibility is a feature and an attack surface at once.}

  \vspace{0.3em}
  {\scriptsize\textcolor{softgray}{Downgrade behavior against 5G devices is
   demonstrated in the research literature. Treat it as a known capability,
   not as a claim about any particular network.}\par}
\end{frame}

\begin{frame}{What would it take to notice?}
  \vspace{0.4em}
  \small
  \begin{itemize}\itemsep6pt
    \item A real tower and a fake one look similar from one measurement.
    \item The signal that matters is a \textbf{pattern}: odd parameters, a
      cell that appears from nowhere, a push to an older generation.
    \item Seeing a pattern means watching many cells at once, over time.
    \item And the pattern changes faster than the standards do.
  \end{itemize}

  \vspace{0.8em}
  \ask{Where in the network could that detector even live?}

  \vspace{0.4em}
  {\small\textcolor{softgray}{Hold that question. It is the one the next section
   was built to answer.}\par}
\end{frame}

\begin{frame}{Why this belongs in an architecture lecture}
  \vspace{0.6em}
  \large
  Every fix in this section was a change to \textbf{who talks to whom, in what
  order, and who is allowed to decide}.

  \vspace{1.0em}
  \normalsize
  That is not a security topic bolted onto an architecture topic.

  \vspace{0.5em}
  It is the same topic.

  \vspace{1.0em}
  \takeaway{Architecture decides what an attacker is allowed to try.}
\end{frame}
